% PROJECT-SPECIFIC SUPPORTING STATEMENT TEMPLATE
% Upload to Overleaf as main.tex. Compile with pdfLaTeX.
% Replace all bracketed text with your own prose.
% Comments beginning with % will not appear in the PDF.
% Follow any word/page limit or prompt given in the application portal.
% The paragraph lengths below are optional drafting guides, not requirements.
% Use connected paragraphs; the prompts are not submission-ready prose.
%
% Project advert (not re-checked from this session):
% https://www.findaphd.com/phds/project/?p198820

\documentclass[11pt,a4paper]{article}
\usepackage[margin=1in]{geometry}
\usepackage[T1]{fontenc}
\usepackage[utf8]{inputenc}
\usepackage{lmodern}
\usepackage{microtype}
\usepackage{setspace}
\usepackage{fancyhdr}
\usepackage[hidelinks]{hyperref}

% ---------- Application details ----------
\newcommand{\ApplicantName}{Bahar Barghbani}
\newcommand{\UniversityName}{The University of Manchester}
\newcommand{\ProjectTitle}{AI-driven Hardware-Software Co-optimisation for
Sustainable and Efficient Computing}
\newcommand{\Supervisors}{Dr Pavlos Petoumenos and Dr Pierre Olivier}

% ---------- Formatting ----------
\setstretch{1.08}
\setlength{\parindent}{0pt}
\setlength{\parskip}{0.7em}
\setlength{\emergencystretch}{2em}
\pagestyle{fancy}
\fancyhf{}
\renewcommand{\headrulewidth}{0pt}
\fancyfoot[C]{\small\thepage}
\hypersetup{pdftitle={Supporting Statement},pdfauthor={Bahar Barghbani}}

\begin{document}

\begin{center}
    {\Large\bfseries Supporting Statement}\par
    \vspace{0.35em}
    {\normalsize\bfseries\ApplicantName}\par
    \vspace{0.25em}
    {\small\UniversityName}\par
    \vspace{0.5em}
    \begin{minipage}{0.95\textwidth}
        \centering\small\ProjectTitle
    \end{minipage}\par
    \vspace{0.35em}
    {\small Supervisors: \Supervisors}
\end{center}
\vspace{0.3em}

% 1. MOTIVATION: approximately 100--130 words.
% Explain your own intellectual reason for choosing this project.
% Connect your interest in memory behaviour and performance to the question
% of what information from software could guide hardware decisions.
% Introduce your current degree briefly; avoid a long autobiographical opening.


I got interested in systems level research in my third semester at Sharif University of Technology. 
During my studies, I found a strong interest in understanding what is happening 
beneath the software level and how different components of a computer interact and cooperate. Moreover, my 
background in physics affects this interest in architecture level research, because I 
feel closer and more connected to it and I find it more physical.
I am particularly interested in memory-focused research and analyzing different mechanisms and innovations 
in the memory hierarchy to improve the performance of the processor. Furthermore, I also enjoy 
making use of software driven information to help with the functionality of the hardware. These reasons together are 
why I am interested in this particular project.

During Computer Structure and Languages course in that semester, I
became familiar with ISAs and assembly programming. I realized I enjoy understanding the operations 
beneath the software level. After that and with passing Computer Architecture and 
Operating Systems courses, I found myself very interested in analyzing 
the operations and interactions between different components of the computer.
Therefore, I chose to deepen my understanding by conducting research in the area and I 
joined Professor Alberto Ros's research at the University of Murcia.

I started working with Prof. Ros since last winter, and I spent 
some time studying the literature and working on reproducing the 
results of an already published paper. However, I realized there are some discrepancies between 
the results and the ones reported in the published paper. After proposing several causes for this issue 
and collaborating with the PhD student who had worked on the project, we decided to move on and use 
an extended version of ChampSim, which was developed to include the Entangling Prefetcher.
Currently, I am working on CVP workload traces to investigate L2 cache placement 
for prefetched instructions. Right now, I am trying different 
policies such as different placement of block heads and block consecutive instructions, 
making use of the confidence of the entangled pair to decide on the pair's placement and working to add L2 prefetching to the overall mechanism and testing 
several settings to figure out a better way that works for these traces 
and outperforms a setting in which all prefetched instructions fill 
L1I cache.

More about my background, I am a bronze medalist in National Physics Olympiad. 
I believe my background in physics was the root of my interest in 
low-level computer research, because I have a physical sense of what 
I am doing and what is happening in the result. After starting my work with 
Professor Alberto and spending a six-week on-site internship with 
Professor \"Ozcan \"Ozt\"urk, which involved implementing compute-intensive and 
memory intensive benchmarks from scratch, as well as studying the state-of-the-art 
methods of using PMU for performance analysis [name of the papers here as well as changing this claim accordingly].
% TODO (Bahar): the sentence starting "After starting my work with Professor
% Alberto..." is incomplete -- it needs a main clause (what happened after?).

% Working on developing a specific yet portable kernel sounds exciting to me, because 
% it involves an intersection of my research interests and background.
% Changing the kernel and tackling with paging and address translation is what stands out to me in this 
% project, because it enables me to work on a broader 
% [Explain why you want to pursue this particular PhD. Identify the question at
% the hardware/software interface that interests you, and describe the experience
% that led you towards it.]
My main motivation to pursue a PhD is to dive deeper into working 
in memory focused research, focusing on cache and front-end of the processor.
In this project in particular, I am most excited about automating 
prefetching using LLM and using AI. Working in the intersection of Machine Learning and 
Memory Optimization is a field I can imagine myself working committedly in.


% 2. ACADEMIC AND OS PREPARATION: approximately 90--120 words.
% Describe your Computer Engineering degree at Sharif and, if useful, your
% major GPA (19.13/20); verify that it is current before submitting.
% Select relevant courses and actual work rather than listing every course.
% Explain what your Operating Systems TA work involved using concrete duties.
% Discuss virtual memory or process mechanisms only if you actually taught,
% implemented, or studied them; TA experience does not imply kernel expertise.
% [Explain your academic preparation for systems research. Give one concrete
% example from relevant coursework, a project, or your Operating Systems teaching
% assistantship that demonstrates your understanding and ability to communicate
% technical ideas. Connect that example to your readiness for this PhD.]
Through my Computer Architecture course, I became familiar with memory and cache and learnt 
some of the basic concepts of memory hierarchy. After I was enrolled in the research in 
Entangling Prefetcher, I developed a strong interest in this area of research and improving 
processor performance through memory optimization.



% 3. RESEARCH WITH ALBERTO ROS: approximately 130--170 words.
% Cover the research question, your individual contribution, experimental
% method, and one lesson that changed your understanding.
% Relevant material may include instruction-cache prefetching, ChampSim,
% Google Sierra workloads, and comparisons of architectural configurations.
% Explain how you distinguish a proposed optimisation from a true bottleneck.
% Verify the current scope and status. Do not imply publication or a successful
% optimisation if the investigation is ongoing. Quantify only verified results.
% [Describe your research with Prof. Alberto Ros. Explain what you are trying to
% understand, what you personally implemented or evaluated, and how you compare
% configurations. Discuss a specific finding, difficulty, or methodological lesson.
% Show how this experience developed your ability to investigate system behaviour
% and assess whether an intervention addresses a meaningful bottleneck.]

% 4. PLANNED WORK WITH OZCAN OZTURK: optional, approximately 40--70 words.
% The research has not properly started. Give it only one or two sentences,
% or omit it if it does not add useful evidence of your direction.
% State the planned topic and why it interests you. Mention any preparation
% only if you have actually completed it, clearly separating it from the plan.
% Do not claim completed experiments, findings, or skills acquired through
% this project. Planned work is motivation, not evidence of research achievement.
% Let your actual contributions with Alberto Ros carry the research argument.
% [If useful, briefly describe the planned performance-analysis project with
% Prof. \"Ozcan \"Ozt\"urk and its intended focus. Make its early status explicit.
% Explain how the topic connects to your developing research interests. Include
% only preparation you have actually completed; otherwise omit preparation claims.]
The research with Prof. \"Ozt\"urk is supposed to start more seriously and beyond preparation
stages. The research question is whether we can come up with useful metrics derived by 
engineering gathered data. The named data is supposed to be gathered from PMU registers.
% 5. PROJECT FIT: approximately 160--210 words.
% Choose one or two aspects you can explain clearly in your own words.
% The advert: LLMs extract compact information from program source code,
% embedded in binaries, to guide hardware prefetching and cache replacement,
% complementing runtime hardware learning at lower implementation cost.
% Explain WHY an aspect interests you; avoid paraphrasing the whole advert.
% Distinguish your transferable skills from knowledge you want to acquire
% (ML, LLMs and compiler/binary techniques are learning goals).
% Potential questions below are prompts for interest, not claimed research gaps:
% - Is prediction confidence enough to choose a cache level, or is another
%   signal needed? Could program-level information provide it?
% - Which hints could reduce the hardware state needed to learn at runtime?
% Connect your experience to a plausible contribution, not a promised solution.
% [Identify one or two aspects of the advertised research that particularly
% interest you. Explain the design trade-off or question you would like to
% understand and how your prior work prepares you to investigate it. Describe
% the ML or compiler knowledge you want to develop,
% and why working with this supervisory team would support that development.]

% 6. GOALS AND CLOSING: approximately 70--100 words.
% Give an honest longer-term direction. Close on your preparation and purpose.
% Avoid generic university rankings, unsupported expertise, or repeating the CV.
% [Explain your longer-term research or professional goals and how this PhD
% would help you pursue them. Close by connecting your experimental preparation,
% interest in systems design, and willingness to learn to the work of the project.]

% BEFORE SUBMISSION:
% - Remove every bracketed prompt; check all personal facts and project claims.
% - Keep the full project title and supervisors identifiable.
% - Check the portal's actual word/page limit and document instructions.
% - Use \%, \&, \_, and \# for literal LaTeX special characters.

\end{document}
